\chapter{Segmento de Clientes}
Para definir con precisión el público objetivo del sistema, se elaboró el Perfil del Cliente basado 
en investigación de campo: entrevistas al Analista de Selección de Tambo, observación directa en las 
oficinas de Recursos Humanos, y encuestas a 20 postulantes reales. El segmento primario es el área de 
Atracción de Talento de Tambo, compuesta por 3 analistas de selección que gestionan en promedio 80 vacantes 
mensuales para perfiles operativos (operarios de tienda, cajeros y reponedores).

\section{ Customer Jobs (Trabajos del cliente)}
\begin{itemize}
  \item \textbf{Funcionales:} Publicar vacantes masivas para operarios de tienda, filtrar perfiles que cumplan con 
  los requisitos mínimos de residencia y disponibilidad, y coordinar entrevistas con Administradores de Tienda 
  en múltiples distritos simultáneamente.

  \item \textbf{Sociales:} Mantener la reputación de Tambo como marca empleadora atractiva, mediante comunicaciones profesionales, 
  oportunas y personalizadas hacia los postulantes durante todo el proceso.

  \item \textbf{Emocionales:} Experimentar seguridad y control al cumplir las cuotas semanales de contratación 
  (promedio: 20 contrataciones/semana) necesarias para las aperturas de nuevas tiendas, sin incurrir en retrasos 
  que generen penalidades operativas.
\end{itemize}

\section{Pains (Dolores)}
\begin{itemize}
  \item \textbf{Frustraciones:} Cada analista invierte entre 4 y 6 horas diarias en la revisión manual de currículos no 
  alineados con el perfil, recibidos por correo electrónico sin estructura estandarizada.
  \item \textbf{Obstáculos:} Cada tienda mantiene sus propios registros de candidatos en hojas de cálculo locales, 
  impidiendo que la Sede Central visualice el pipeline de contratación en tiempo real.
  \item \textbf{Riesgos:} El $35\%$ de los candidatos aptos abandona el proceso antes de ser contactado para entrevista, debido a demoras superiores a 5 días entre la postulación y la primera comunicación.
\end{itemize}

\section{Gains (Ganancias)}
\begin{itemize}
  \item \textbf{Resultados esperado (cuantificado):} Reducir el Tiempo de contratación de 15-20 días a 5-7 días hábiles, 
  disminuyendo la tasa de abandono de candidatos aptos del $35\%$ al $10\%$.
  \item \textbf{Resultados deseado:} Disponer de una base de datos centralizada y segmentada geográficamente, 
  que permita reutilizar candidatos descartados en convocatorias futuras para tiendas en la misma zona.
  \item \textbf{Resultado inesperado de valor:} Generación automática de reportes de métricas del proceso
   (candidatos por etapa, tasa de conversión, tiempo por fase) para la toma de decisiones estratégicas en RRHH.
\end{itemize}
